Er zijn verschillende manieren om een regel te verwijderen. Je kan de regel zoals je die hebt ingevoerd voorzien van een \textbf{delete} argument inplaats van \textbf{allow} of \textbf{deny}:
\begin{lstlisting}[language=bash]
sudo ufw delete snmp-trap
\end{lstlisting}

Een makkelijkere manier kan zijn door eerst de status van UFW op te vragen met nummers en dan de regel te deleten op basis van zijn regelnummer:
\begin{lstlisting}[language=bash]
sudo ufw status numbered
\end{lstlisting}
Levert een autoput op die lijkt op:
\begin{lstlisting}[language=bash]
Status: active

     To           Action     From
     --           ------     ----
[ 1] 42/tcp       ALLOW IN   Anywhere
[ 2] 42/udp       DENY IN    Anywhere
[ 3] 42/tcp (v6)  ALLOW IN   Anywhere
[ 4] 42/udp (v6)  DENY IN    Anywhere
\end{lstlisting}
Op basis van bovenstaande kunnen we bijvoorbeeld regel 2 verwijderen. We verwijderen dan ook alleen regel 2, regel 4 (de IPv6 variant) blijft bestaan.
\begin{lstlisting}[language=bash]
sudo ufw delete 2
\end{lstlisting}
Verwijder regel 2 en controlleer dit met een genummerde status.

